\scanpage{4}

Ceci dit, si on se donne un
multimarquage de $\partial X$, savoir pour
chaque $i\in I$ un $n_i\in\mathbb N^*$ et un
sous-\inserted{$\mathbb U_{n_i}$-}torseur
$T_i(n_i)\subset T_i$,
\deleted{\unsure{les groupes}} i.e.
un $s_i\in T_i^{\wedge n_i}$, d'où
$\deleted{\unsure{\uline{s}\in}}$
\[
\uline{s}\in\prod_{i\in I}T_i^{\wedge n_i},
\]
alors le groupe de Teichmüller \unsure{multimarqué}
$S_{\unsure{*|1}}T(\unsure{X},\uline{s})$,
généralisant le précédent, est
plongé dans $\operatorname{SRT}(X)$ de façon naturelle,
en prenant un $s_i^0$ dans chaque $T_i(n_i)$, d'où
$\uline{s}^0\in\prod T_i$ et
$ST(X,\uline{s}^0)\subset\operatorname{SRT}(X)$,
et on prend le ss-groupe
\[
ST(X,\uline{s})=ST(X,\uline{s}^0).
\prod_{i\in I}\frac1{n_i}\mathbb Z_i
\]
où $\mathbb Z_i$ est le $i$ème facteur de
$\mathbb Z^I\subset\operatorname{SRT}(X)$.

\margin{3}
